\section{Introduction and conclusions}

\subsection{Semigroups and equational laws}

% TODO

magmas, semigroups, equational theories,

Should we have a section setting up Mathematical Foundations more precisely as we do in the ETP?

Note about the numbering we use, which is strange from the semigroup point of view but which comes from the ETP.

\subsection{Outcome}

The finite and infinite implication graphs are identical.

\begin{theorem}
  Any associative equational law of order at most~$4$ is equivalent to one of $122$ laws.
  Any conjunction of such laws sits in one of $456$ equivalence classes, given in \autoref{app:alltheories} by maximal conjunctions of equations.
  Implication and conjunction are the subset and join relations.
\end{theorem}

\subsection{Automation vs formalization}

Distinguish two axes: manual--automated and informal--formalized, so
\begin{itemize}
\item manual \& informal is what we do chatting on Zulip
\item manual \& formalized is writing one theorem (such as the end-to-end theorem)
\item automated \& informal is the use of ATPs such as Mace4/Prover9 to initially get the implication graph
\item automated \& formal is the Vampire--Lean integration
\end{itemize}

\subsection{Future work}

Extend to order~$5$ and then~$6$ to see if there ends up being any distinctions between finite and infinite implication graphs.
Actually, maybe there are theorems forbidding this in associative equations, who knows?

Much harder: all conjunctions of magma equations of order up to $4$.  This will at last include the case of groups, Boolean algebras, etc.
